\section{Jagtetiske regler}

\begin{frame}
  \vspace{3em}
  \begin{center}
    {\LARGE \textbf{Jagtetiske regler}}
    \\
    \vspace{1em}
    Danmarks Jægerforbund, Naturstyrelsen, Vildtforvaltningsrådet
  \end{center}
\end{frame}

\begin{frame}{Generelle hensyn – Jagtens omdømme \& etik}
\begin{itemize}
\item Alt jagt skal følge lovgivning samt jagtetiske regler
\item Jagt er mere end skydning – anvend nedlagt vildt mest muligt (føde, pels, trofæ)
\item Vildt behandles med respekt – også uanvendeligt vildt (nedgravning/destruktion)
\item Efterlad naturen uden spor – tag alt affald med
\item Undgå jagt på udsat vildt med unaturlig adfærd
\item Kun nødvendig fodring – ænder primært korn/frø på land
\item Forsøg at være et forbillede – lær af erfarne jægere
\end{itemize}
\end{frame}

\begin{frame}{Hensyn til vildtbestandene}
\begin{itemize}
\item Jag kun en del af det naturlige overskud – brug forsigtighedsprincippet
\item Kend dit vildt – skyd kun når du er sikker på arten
\item Tilpas jagttryk til terrænets bæreevne – ikke til leje eller antal jægere
\item Trækkende vildt = fælles ansvar – skån hvileområder og fouragering
\item Begræns forstyrrelser: max 3–4 klap-/drivjagt pr. sæson pr. areal
\item Minimum 2–3 ugers interval mellem jagter på samme lokalitet
\item Afbalanceret afskydning (køn, alder) – forvalt som du ønsker naboen gør det
\end{itemize}
\end{frame}

\begin{frame}{Hensyn til individet – Effektiv aflivning}
\begin{itemize}
\item Brug korrekt våben/ammunition – indskyd og træn regelmæssigt
\item Respektér dine skydebegrænsninger – træn på banen, ikke på jagt!
\item Maks. afstande hagljagt (forsigtige anbefalinger):
  \begin{itemize}
  \item Råvildt 20 m \quad Gæs/ræv 25 m \quad Øvrigt 30 m
  \end{itemize}
\item Riffel: Fast anlæg – max 100 m – frit skudfelt – ingen hals/hovedskud
\item Kun ét mål ad gangen – ingen skift af dyr ved anskydning
\item Saml alt op – eftersøg ihærdigt – hjælp hinanden
\end{itemize}
\end{frame}

\begin{frame}{Praktiske skudregler \& eftersøgning}
\begin{itemize}
\item Hagl: max $\sim$3 patroner pr. flyvende fugl, $\sim$2 pr. løbende vildt
\item Riffel: skudantal $\approx$ antal nedlagt vildt
\item Hold øje med anskudt vildt – marker præcist – instruer hundefører
\item Havjagt: forfølg dykkende anskudte vedholdende – stop jagt ved vanskelige forhold
\item Affang anskudt vildt hurtigt og dyreværnsmæssigt korrekt
\item Aldrig dublé råvildt med hagl – begræns dublé i øvrigt
\end{itemize}
\end{frame}

\begin{frame}{Hensyn til andre naturbrugere}
\begin{itemize}
\item Vær venlig og vejled ift. jagtaktivitet – irettesæt kun når nødvendigt
\item Undgå unødvendig lukning af skov – brug vejledende skilte
\item Stor forsigtighed med haglnedfald nær andre mennesker
\item Jag ikke hvor det er til klar ulempe for andre
\item Strand-/havjagt: hold god afstand til alle andre brugere
\end{itemize}
\end{frame}

\begin{frame}{Hensyn til medjægere \& nabosamarbejde}
\begin{itemize}
\item Respektér andres jagtret – samarbejd med nabojægere
\item Aftal jagthyppighed, afskydning og eftersøgning
\item Skyd aldrig vildt der drives fra naboterræn
\item Påtal ulovlig eller uetisk jagt – altid
\item Vis storsind ved tvivl om vildtets tilhørsforhold
\item Sikker våbenhåndtering – altid
\end{itemize}
\end{frame}

\begin{frame}{Sikkerhed}
\begin{itemize}
\item Røde hattebånd / høj synlighed – sikkerhedsbriller ved hagl
\item Våben uladt, knækket/åben, lodret mellem poster
\item Minimum 40$^{\circ}$ sikkerhedsvinkel – ekstra forsigtighed ved rikochet
\item Ingen lav skydning mod uoverskuelig baggrund
\item Riffel: altid sikkert kuglefang
\item Strand-/havjagt: moderne sikkerhedsudstyr – ingen unødige chancer
\item Følg Søsportens Sikkerhedsråds anbefalinger
\end{itemize}
\end{frame}